\subsection{Convenciones}

\subsubsection{Regla: variable de excepción nombrada \texttt{ex}}

\textbf{Regla:} Toda variable de excepción que sí se utiliza en el bloque \texttt{catch} se nombra \texttt{ex}.

\textbf{Descripción:} Convención propia del proyecto para mantener los bloques \texttt{catch} predecibles en todo el equipo (McConnell, 2004).

\textbf{Código correcto:}
\begin{lstlisting}[style=csharpstyle]
try
{
    _board.MovePlayer(player, destination);
}
catch (InvalidMoveException ex)
{
    _logger.Warning($"Invalid move attempted by {player.Name}: {ex.Message}");
}
\end{lstlisting}

\textbf{Código incorrecto:}
\begin{lstlisting}[style=csharpstyle]
try
{
    _board.MovePlayer(player, destination);
}
catch (InvalidMoveException exception)
{
    _logger.Warning($"Invalid move attempted by {player.Name}: {exception.Message}");
}
\end{lstlisting}

\subsubsection{Regla: orden de las cláusulas \texttt{catch}}

\textbf{Regla:} Varias cláusulas \texttt{catch} en un mismo \texttt{try} se ordenan del tipo más derivado al menos derivado (Microsoft, 2025a).

\textbf{Descripción:} Con herencia directa entre los tipos, como abajo, el compilador rechaza el orden incorrecto con un error: no es solo preferencia de estilo.

\textbf{Código correcto:}
\begin{lstlisting}[style=csharpstyle]
try
{
    _board.MovePlayer(player, ParsePosition(rawDestination));
}
catch (ArgumentOutOfRangeException ex)
{
    _logger.Warning($"Position out of board bounds for {player.Name}: {ex.Message}");
}
catch (ArgumentException ex)
{
    _logger.Warning($"Invalid position format for {player.Name}: {ex.Message}");
}
\end{lstlisting}

\textbf{Código incorrecto (no compila --- CS0160):}
\begin{lstlisting}[style=csharpstyle]
try
{
    _board.MovePlayer(player, ParsePosition(rawDestination));
}
catch (ArgumentException ex)
{
    _logger.Warning($"Invalid position format for {player.Name}: {ex.Message}");
}
catch (ArgumentOutOfRangeException ex)
{
    _logger.Warning($"Position out of board bounds for {player.Name}: {ex.Message}");
}
\end{lstlisting}
